\section{Verification Environment}
\label{sec:verification}

Ensuring the functional correctness of the RTL before logic synthesis and Physical Design (APR) was an absolute paramount. To do this, we architected a robust, modular Universal Verification Methodology (UVM) testbench driving a native 100 MHz simulation environment.

\subsection{UVM Testbench Architecture}
The environment leverages standard UVM phasing to stimulate the NPU's RTL under real-world constraints:
\begin{itemize}
    \item \textbf{Sequence \& Driver:} The \texttt{my\_uvm\_sequence} class is responsible for parsing the exact same hexadecimal test files generated by the Python pre-processing phase. It actively drives the \texttt{NN\_NUM\_INPUTS} (784) pixel sequences into the Device Under Test (DUT) via the \texttt{nn\_if} SystemVerilog interface, asserting the \texttt{wr\_en} signal cycle-by-cycle exactly as an external SRAM or image sensor would stream data to the chip.
    \item \textbf{Monitor \& Scoreboard:} Once the hardware calculates the final output and asserts the \texttt{inference\_done} flag, the \texttt{my\_uvm\_monitor} captures the \texttt{predicted\_class} alongside its 16-bit confidence \texttt{score}. These values are broadcasted via a UVM Analysis Port to the \texttt{my\_uvm\_scoreboard}. The scoreboard independently reads the ``ground truth'' labels exported by PyTorch and strictly compares the RTL output against the expected classification.
\end{itemize}

\subsection{Functional Coverage \& Coverage Hooks}
Simulating the RTL against input vectors is not sufficient proof of a complete verification cycle without Functional Coverage metrics. We extensively utilized SystemVerilog \texttt{covergroup} structures to gather runtime statistics:
\begin{itemize}
    \item \textbf{Prediction Coverage (\texttt{cg\_prediction}):} We defined independent \texttt{coverpoints} for \texttt{predicted\_class} and \texttt{expected\_class} spanning bins [0:9]. A \texttt{cross} coverage block mathematically tracked testing sequences to verify we stimulated and correctly classified every possible digit from 0 to 9 without any hardcoded blind spots.
    \item \textbf{Layer Activations Coverage (\texttt{cg\_layer\_activations}):} This was our most technically rigorous check. By embedding internal probes deep within the RTL logic (connecting to \texttt{l0\_out} through \texttt{l3\_out}), our covergroups tracked whether every single neuron's ReLU output toggled effectively (firing into the \texttt{positive} bin) or remained dormant (\texttt{zero} bin). 
\end{itemize}

\subsection{UVM Verification Results}
The final \texttt{make uvm\_sim} regression produced an exemplary clean report as shown in Fig.~\ref{fig:uvm_summary}. The UVM Report Summary recorded exactly \textbf{22 UVM\_INFO} messages, \textbf{0 UVM\_WARNING}, \textbf{0 UVM\_ERROR}, and \textbf{0 UVM\_FATAL} violations, confirming a fully clean verification run. The report counts by component ID show that all verification agents operated correctly: the Driver (\texttt{[DRV]}) issued 4 transactions, the Monitor (\texttt{[MON]}) captured 7 observation events, the Scoreboard (\texttt{[SCB]}) performed 5 comparison checks, and the Sequencer (\texttt{[SEQ]}) processed 2 sequence items.

\begin{figure}[ht]
    \centering
    \includegraphics[width=0.85\columnwidth]{Fig/uvm_report_summary.png}
    \caption{UVM Report Summary showing zero errors and zero fatal violations across all verification components.}
    \label{fig:uvm_summary}
\end{figure}

\subsection{Performance Metrics}
Fig.~\ref{fig:uvm_results} presents the detailed simulation output of a single UVM test case. The Monitor captures the inference completion event with a predicted class of 9 and a maximum confidence score of 11633, which the Scoreboard confirms as a \textbf{TEST PASSED} result by comparing against the expected label 9 from the golden reference file.

Critically, the embedded performance counters reveal the NPU's temporal behavior:
\begin{itemize}
    \item \textbf{First Write Cycle:} 6 (the testbench begins streaming data at cycle 6 after reset release).
    \item \textbf{Inference Done Cycle:} 1745 (the hardware asserts \texttt{inference\_done} at cycle 1745).
    \item \textbf{Total Latency:} 1739 clock cycles from data input to classification output.
    \item \textbf{@100MHz Inference Time:} \textbf{17.39 $\mu$s}, demonstrating the NPU achieves real-time classification performance suitable for high-speed edge deployment.
\end{itemize}

\begin{figure}[ht]
    \centering
    \includegraphics[width=\columnwidth]{Fig/uvm_test_results.png}
    \caption{UVM single-class test output showing correct inference (Predicted: 9, Expected: 9) and performance summary with 1739-cycle total latency.}
    \label{fig:uvm_results}
\end{figure}

\subsection{Batch Simulation: All 10 MNIST Classes}
To comprehensively validate the NPU across the entire MNIST digit space (0--9), we developed a batch simulation testbench (\texttt{nn\_tb\_all.sv}) that sequentially loads and classifies one representative test image per class. Fig.~\ref{fig:batch_sim} shows the final batch simulation output.

The batch runner iterates through all 10 classes, copying the corresponding \texttt{x\_test\_class\{i\}.txt} and \texttt{y\_test\_class\{i\}.txt} files as DUT stimulus. Each class produces a correct prediction with its associated confidence score. The final summary confirms:
\begin{itemize}
    \item \textbf{Total Tests:} 10
    \item \textbf{Passed:} 10
    \item \textbf{Failed:} 0
    \item \textbf{Result:} \textbf{ALL TESTS PASSED}
\end{itemize}

This 100\% classification accuracy across all 10 digit classes provides strong evidence that the NPU architecture faithfully preserves the trained model's inference capability after fixed-point quantization, and that no systematic hardware bugs exist in the datapath, weight ROM indexing, or FSM control logic.

\begin{figure}[ht]
    \centering
    \includegraphics[width=\columnwidth]{Fig/batch_simulation.png}
    \caption{Batch simulation results showing all 10 MNIST digit classes (0--9) correctly classified with 100\% accuracy. The summary confirms 10/10 tests passed with zero failures.}
    \label{fig:batch_sim}
\end{figure}

\subsection{Waveform Analysis}
Fig.~\ref{fig:waveform} presents a comprehensive waveform capture from the Cadence Xcelium SimVision debugger, illustrating the NPU's internal signal activity across the full batch simulation of all 10 classes. The waveform is organized into five major signal groups:

\begin{itemize}
    \item \textbf{DUT\_IO:} The top-level input/output signals including \texttt{wr\_en}, \texttt{din}, \texttt{in\_full}, \texttt{inference\_done}, \texttt{predicted\_class}, and \texttt{max\_score}. Ten distinct inference pulses on \texttt{inference\_done} are clearly visible, each corresponding to one digit class.
    \item \textbf{RESULTS:} The \texttt{pass\_count} and \texttt{fail\_count} registers, confirming 10 passes accumulating over time with zero failures.
    \item \textbf{FSM:} The finite state machine progression through \texttt{S\_IDLE}, \texttt{S\_LOAD\_IMG\_RUN}, \texttt{S\_MAC\_PASS}, \texttt{S\_WAIT\_MAC}, \texttt{S\_DRAIN}, and \texttt{S\_DONE} states. The repetitive FSM pattern across all 10 test cases validates deterministic control flow. The \texttt{img\_cnt} signal reaches 784 during each image load phase, confirming all input pixels are consumed.
    \item \textbf{FIFO:} The \texttt{fifo\_empty} and \texttt{fifo\_rd\_en} signals demonstrate correct handshaking between the testbench data source and the NPU's internal input buffer.
    \item \textbf{MAC\_CTRL:} The \texttt{mac\_start\_in}, \texttt{mac\_valid\_in}, and \texttt{mac\_last\_in} control signals orchestrating the 32 parallel MAC lanes across the 4 network layers.
\end{itemize}

\begin{figure*}[ht]
    \centering
    \includegraphics[width=\textwidth]{Fig/waveform.png}
    \caption{Full waveform capture of the batch simulation across all 10 MNIST classes. Signal groups include DUT I/O, test results, FSM state transitions, FIFO control, MAC array orchestration, and Argmax output.}
    \label{fig:waveform}
\end{figure*}
